\subsection*{S11. Key names that are not OCaml identifiers}

JSON permits any string as a member name; OCaml record fields must be
lowercase identifiers drawn from a restricted ASCII set and must not collide
with keywords. Several distinct failures arise:

\begin{lstlisting}[style=json]
{ "my-field": 1,      -- hyphen is illegal in an identifier
  "type": 2,          -- an OCaml keyword
  "1abc": 3,          -- leading digit
  "Name": 4,          -- leading uppercase; record fields must start lowercase
  "caf\u00e9": 5,   -- non-ASCII; OCaml identifiers are ASCII only
  "": 6 }             -- empty
\end{lstlisting}

The capitalized case is the one most easily overlooked: a capitalized key is
entirely ordinary JSON, emitted by most .NET and Java producers, and cannot be
a record field name at all.

The binding constraint is \textbf{injectivity}. Two distinct JSON keys must
never map to the same OCaml name, or two columns silently become one. Naive
replacement fails at once: \texttt{my-field} and \texttt{my\_field} both
become \texttt{my\_field}.

\options
\begin{enumerate}[label=\textbf{(\alph*)}]
  \item \textbf{Unconditional injective mangling.} A total scheme applied to
        every key, with the original key retained in the schema. Names are
        stable and always distinct, and frequently ugly. \selected
  \item \textbf{Reject} any key that is not already a legal OCaml identifier.
  \item \textbf{Context-sensitive mangling} --- use the readable form and
        escape only when that particular object would otherwise collide.
  \item \textbf{Naive mangling with collision detection} --- simple
        replacement, erroring when two keys collide.
\end{enumerate}

\decision{Unconditional injective mangling. Every key is mangled by the same
total function, whether or not it is already legal, and the original key is
recorded in the schema alongside the mangled name.}

\implnote{A concrete scheme, sufficient to implement and adjustable in its
cosmetic details so long as injectivity is preserved:
\begin{enumerate}[label=\arabic*.]
  \item Process the key as UTF-8 bytes. A byte in \texttt{a-z}, \texttt{A-Z},
        \texttt{0-9} or \texttt{'} passes through unchanged; every other byte,
        \emph{including} \texttt{\_}, becomes \texttt{\_} followed by two
        lowercase hexadecimal digits. Escaping \texttt{\_} is what makes
        \texttt{\_} an unambiguous escape introducer.
  \item If the result is empty or does not begin with a lowercase letter,
        prefix it with \texttt{f\_}.
  \item If the result is an OCaml keyword, append \texttt{\_}.
\end{enumerate}
Under this scheme \texttt{my-field} becomes \texttt{my\_2dfield},
\texttt{type} becomes \texttt{type\_}, \texttt{1abc} becomes
\texttt{f\_1abc}, and \texttt{Name} becomes \texttt{f\_Name}.

That this scheme is injective is a property to be checked rather than assumed;
it is exactly the kind of small lemma worth proving once the formal
development resumes.}

\rationale{Option (b) rejects a large fraction of real-world JSON: hyphenated
and capitalized keys are commonplace, and the user usually cannot rename the
source data. Option (c) makes a generated name depend on what else appears in
the same object, so adding a key to the corpus can rename an existing field
and break code written against the generated module --- the same
data-dependence problem as S9, in a milder form, since an object's key set is
structure rather than data.

Reconstruction does not require the mangling to be invertible by inspection.
The original key is stored in the schema, so injectivity alone is what is
needed: it prevents two columns from merging, and the stored key restores the
exact name on output.}

\limitation{Mangled names are less readable than the keys they came from, and
the delivery vehicle sharpens the cost: these modules are read by people in an
editor rather than consumed as a build artifact. A middle path was considered
--- mangle readably but reserve the escaped forms, so collisions are
impossible by construction --- which recovers most of (c)'s readability while
retaining (a)'s stability. It remains available as a later refinement.

Module names, which must instead begin with an uppercase letter, need a
separate rule; that belongs to S22 along with collisions between generated
names.}
